\section{Decision-time search}
\label{sec:search}

The learned actor's prior looks one turn ahead against a handful of guessed
opponent moves. The search agent replaces that guess with the official
simulator. It samples complete hypothetical opposing teams, plays every
pairing of our and the opponent's joint actions for one turn, and solves the
resulting matrix game. Figure~\ref{fig:search} shows the procedure.

\begin{figure}[htbp]
\centering
\resizebox{\linewidth}{!}{%
\begin{tikzpicture}[x=1cm,y=1cm]
  \node[box,text width=34mm] (trk) at (0,0) {Public tracker\\HP, status, boosts,\\revealed sets, field};
  \node[box,text width=34mm] (use) at (0,-2.2) {Usage prior\\Smogon chaos table,\\September 2026};
  \node[warmbox,text width=27mm] (smp) at (4.3,-1.1) {Sample $K=6$\\worlds $w$\\(our sets exact)};
  \node[note,anchor=north] at (4.3,-1.95) {3 Node engines,\\2 worlds each};
  \foreach \k in {3,2,1}{
    \draw[draw=ink,fill=pale,rounded corners=2pt] (6.6+0.12*\k,-0.2-0.12*\k) rectangle (10.4+0.12*\k,-2.5-0.12*\k);
  }
  \node[font=\scriptsize,anchor=south] at (8.6,-0.15) {world $w$: a rebuilt simulator \code{Battle}};
  \foreach \i in {0,...,5}{ \foreach \j in {0,...,6}{
     \pgfmathsetmacro{\shade}{mod(7*\i+3*\j,9)*8+10}
     \fill[seriesA!\shade!white] (7.7+0.3*\j,-0.95-0.22*\i) rectangle (8.0+0.3*\j,-1.17-0.22*\i);
  }}
  \node[note] at (8.75,-0.72) {theirs $b$};
  \node[note,rotate=90] at (7.35,-1.6) {ours $a$};
  \node[font=\scriptsize,anchor=north] at (8.6,-2.75) {$M_w[a,b]$: value after one simulated turn};
  \node[warmbox,text width=29mm] (rm) at (13.3,-0.3) {Regret matching\\2000 iterations\\$\to(\bar x_w,\bar y_w)$};
  \node[warmbox,text width=29mm] (bl) at (13.3,-2.3) {Opponent model\\$q_w=\tfrac12\bar y_w+\tfrac12g_w$};
  \node[box,text width=58mm] (avg) at (7.6,-4.3) {$\displaystyle a^*=\arg\max_a\frac1{|W_a|}\sum_{w\in W_a}M_w[a,\cdot]\,q_w$};
  \node[box,text width=26mm] (out) at (13.3,-4.3) {click $a^*$\\median 2.3\,s};
  \draw[arr] (trk) -- (smp); \draw[arr] (use) -- (smp);
  \draw[arr] (smp) -- (6.65,-1.1);
  \draw[arr] (10.9,-0.6) -- (rm.west);
  \draw[arr] (rm) -- (bl);
  \draw[arr] (bl.south) -- ++(0,-0.5) -| ([xshift=10mm]avg.north);
  \draw[arr] (avg) -- (out);
\end{tikzpicture}}
\caption{The search agent. Each world fixes complete hypothetical opposing
sets and copies all public state. Screening keeps at most 24 of our joint
actions and 20 of the opponent's. The payoff matrix of a world therefore needs
at most 480 one-turn simulations, plus up to $4|\B|+4|\A|$ for screening. Team preview is delegated to the learned actor.}
\label{fig:search}
\end{figure}

\subsection{Public tracker}
\code{search/tracker.py} reads only our own player stream. For each
Pok\'{e}mon it tracks:
\begin{itemize}
  \item HP fraction (exact HP when the denominator is not 100), status, sleep
    and toxic counters, boosts and volatiles;
  \item revealed moves, item (or its loss) and ability;
  \item Mega state, last move, and the turn it switched in.
\end{itemize}
For the field it tracks weather, terrain, pseudo-weather and side conditions,
each with its start turn. A condition's remaining duration is estimated as
the base duration minus elapsed turns. If that would be non-positive, the
extended duration is assumed. Open team sheets are parsed from
\code{|showteam|} lines.

\subsection{Sampling hidden sets from usage statistics}
\label{sec:prior}
The prior is Smogon's monthly \emph{chaos} usage table for the format
(September 2026, all ratings, about 1.6 million battles)
\citep{smogonstats}. For each usage entry $e$, which is a species or a Mega
forme, the table gives a raw count $c_e$ and normalized frequencies
$P_e(\cdot)$ that sum to one over each category. The categories are moves,
items, abilities and the 60 most common nature-and-stat-point spreads.

\begin{lemma}[Forme posterior under conditional independence]
\label{lem:forme}
Suppose an opposing Pok\'{e}mon's entry is $e$ with prior probability
proportional to $c_e$. Suppose also that, given $e$, the event ``move $m$ is in
the moveset'' has probability $\pi_e(m)$ independently across moves. Then
after observing the revealed set $R$, $\Prob(e\mid R)\propto
c_e\prod_{m\in R}\pi_e(m)$. If, moreover, $\pi_e(m)=c\,P_e(m)$ for a constant
$c$ shared by all entries (with four move slots, $c\approx4$), then
\begin{equation}
  \Prob(e\mid R)\propto c_e\prod_{m\in R}P_e(m).
  \label{eq:forme}
\end{equation}
\end{lemma}
\begin{proof}
The first statement is Bayes' rule with the likelihood factorized by the
independence assumption. Under the proportionality assumption, the product
becomes $c^{|R|}\prod_{m\in R}P_e(m)$. The factor $c^{|R|}$ is the same for every
entry, so it cancels in the normalization.
\end{proof}

The implementation follows Equation~\eqref{eq:forme}, with three adjustments:
\begin{itemize}
  \item unseen moves get a floor of 0.002 on the share scale, so a surprising
    reveal cannot zero out every entry;
  \item an observed non-stone item, or a spent Mega, removes the Mega entries;
  \item an observed Mega keeps only matching Mega entries.
\end{itemize}
Given $e$, the unrevealed moves are drawn without replacement in proportion
to $P_e$. Item, ability and spread are each drawn from their own marginals. A
Mega entry takes its most common item (assumed to be the stone) and the base
forme's abilities. Unseen previewed Pok\'{e}mon fill the remaining slots,
drawn without replacement in proportion to species usage. When open sheets
were exchanged, species, item, ability, moves and nature are copied exactly,
and only stat points are sampled.

\caveat{Lemma~\ref{lem:forme} is exact only under its independence
assumption. Real movesets have four slots, so move inclusions are negatively
dependent. Item, ability and spread are correlated in real sets, and the
\code{Teammates} table is loaded but unused. The Item Clause is not enforced
across a sampled team. The result is a tractable proposal distribution, not a
calibrated posterior over the opponent's team.}

\subsection{Rebuilding a world}
Each world becomes a fresh simulator \code{Battle} built from both packed
teams. Public state is then written onto it:
\begin{itemize}
  \item formes, abilities, items, exact HP for us and $\mathrm{round}(h\cdot
    \mathrm{maxHP})$ for them, statuses with their counters, and our own PP;
  \item a whitelist of volatiles, such as Protect's stall counter, Substitute,
    Taunt, Encore, Disable and choice lock. The agent does not send last
    moves, so the engine's last-move restoration is unused, and the opponent's
    PP is never set;
  \item weather, terrain and side conditions with their estimated durations.
\end{itemize}
If a side has used its Mega, Mega Evolution is disabled for the whole side.
Our side uses exact sets, HP and PP from our request.

\subsection{The one-turn game in each world}
Our candidates are the Python-legal choices $\A(r)$, pruned in two ways.
Damaging moves at our partner are removed, and when a Mega is available only
Mega-containing choices are kept. Each surviving choice is rewritten by move
\emph{identifier}, because a locked request lists only one move and slot
numbers differ in the rebuilt battle. The opponent's set $\B$ comes from the
engine's own enumeration (Section~\ref{sec:actions}).

\paragraph{Screening.} If $|\B|>20$ or $|\A|>24$, the engine takes four
evenly spaced probe actions $P\subset\A$. It keeps the 20 opponent actions
with the lowest mean value against $P$, which are the best for the opponent.
It then keeps the 24 of our actions with the highest mean value against the
first four kept replies.

\paragraph{Payoffs.} For each remaining pair, the engine restores the world
state, gives the cell its own seed, submits both choices, runs exactly one
turn and evaluates the result:
\begin{equation}
  M_w[a,b]=V\big(\mathrm{Turn}(B_w;a,b;\xi_{w,a,b})\big).
\end{equation}
If the simulator rejects any cell, that cell's whole row (our action) or
column (their action) is removed.

\paragraph{Evaluation function.} A finished battle scores $\pm30$, or 0 for a
tie. Otherwise
$V=S(\text{us})-S(\text{them})+\mathrm{TR}$, where
\begin{equation}
  S=\sum_{\text{alive }p}\big[1+1.2h_p-\mathrm{pen}(p)+\mathbf 1_{\text{active}}\,\beta(p)\big]
   +0.10\cdot\text{Tailwind turns}+0.06\sum\text{screen turns}.
\end{equation}
Status penalties are: paralysis 0.30, freeze 0.50, poison 0.12, toxic 0.25,
sleep $0.45\min(1,\text{turns}/2)$, and burn 0.45 if Attack exceeds Special
Attack, else 0.12. The active-Pok\'{e}mon term $\beta$ rewards offensive,
defensive and speed boosts and a Substitute. It penalizes Perish counts,
confusion, Taunt, Encore, Yawn and Leech Seed. TR is the Trick Room term: it
adds $0.10\cdot\text{turns}\cdot2\clip((\bar v_{\text{them}}-\bar v_{\text{us}})/40,-1,1)$,
where $\bar v$ is each side's mean Speed, so slower teams gain value under Trick Room.

\subsection{Solving the matrix game}
The engine runs simultaneous regret matching on the zero-sum game in which we
maximize $M$ \citep{hartmascolell}. With cumulative regrets $R_t$ for us and
$C_t$ for the opponent,
\begin{align}
  x_t(i)&\propto[R_t(i)]^+,\qquad y_t(j)\propto[C_t(j)]^+\quad(\text{uniform if all}\le0),\\
  R_{t+1}(i)&=R_t(i)+(My_t)_i-x_t^\top My_t,\qquad
  C_{t+1}(j)=C_t(j)+x_t^\top My_t-(x_t^\top M)_j,
\end{align}
for $T=2000$ iterations. It returns the average strategies $\bar x$ and $\bar y$.

\begin{proposition}[Average regret gives an approximate equilibrium]
\label{prop:rm}
Let $\varepsilon_1$ and $\varepsilon_2$ be the two players' average external
regrets after $T$ rounds. Then $(\bar x,\bar y)$ is an
$(\varepsilon_1+\varepsilon_2)$-equilibrium:
$\max_i(M\bar y)_i-\min_j(\bar x^\top M)_j\le\varepsilon_1+\varepsilon_2$.
For regret matching with payoff range $\Delta$,
$\varepsilon_1\le\Delta\sqrt{|\A|/T}$ and $\varepsilon_2\le\Delta\sqrt{|\B|/T}$.
\end{proposition}
\begin{proof}
Let $\bar v=T^{-1}\sum_tx_t^\top My_t$. By linearity,
$T^{-1}\sum_t(My_t)_i=(M\bar y)_i$, so our average regret is
$\max_i(M\bar y)_i-\bar v\le\varepsilon_1$. Likewise the opponent's is
$\bar v-\min_j(\bar x^\top M)_j\le\varepsilon_2$. Adding the two gives the
claim. Because
$\min_j(\bar x^\top M)_j\le\bar x^\top M\bar y\le\max_i(M\bar y)_i$, neither
player can gain more than $\varepsilon_1+\varepsilon_2$ by deviating. The
regret rates are the standard bounds for regret matching
\citep{hartmascolell,zinkevich2007}.
\end{proof}

\intuition{With $|\A|=24$, $|\B|=20$ and $T=2000$, the worst-case gap is
$\Delta(\sqrt{24}+\sqrt{20})/\sqrt{2000}\approx0.21\Delta$. When a terminal
$\pm30$ is reachable, $\Delta$ can be 60, which makes the bound loose. In
practice regret matching usually converges much faster than its worst case.
The bound justifies $\bar y$ only as an approximate equilibrium, and only of
this one-turn surrogate game.}

\subsection{The opponent model and the decision rule}
The agent does \emph{not} play $\bar x$. It treats the opponent as a mixture of
two types:
\begin{equation}
  q_w=\alpha\,\bar y_w+(1-\alpha)\,g_w,\qquad
  g_w(b)\propto\exp\!\big(-(c_b-\min_{b'}c_{b'})/\tau_o\big),\qquad
  c_b=\frac1{|\A|}\sum_aM_w[a,b],
\end{equation}
with $\alpha=\tfrac12$ and $\tau_o=1$. One type plays the equilibrium. The
other is a soft best response to our playing uniformly, a level-1 quantal
response in the sense of \citet{mckelvey1995}. Each world then gives our action
values $U_w(a)=(M_wq_w)_a$. Let $W_a$ be the worlds in which $a$ was valid, and
$W_{\mathrm{ok}}$ the worlds that produced any result. The decision is
\begin{equation}
  a^*=\arg\max_a\Big[\frac1{|W_a|}\sum_{w\in W_a}U_w(a)-0.5\,\mathbf 1\{|W_a|<|W_{\mathrm{ok}}|\}\Big].
  \label{eq:decision}
\end{equation}

\begin{proposition}[Guarantee and exploitability of a best response]
\label{prop:br}
Let $v^*$ be the value of a world's matrix game. For any opponent mixture $y$,
a best response $a^*\in\arg\max_a(My)_a$ satisfies $(My)_{a^*}\ge v^*$. Its
worst case, $\min_bM[a^*,b]$, can still be strictly below $v^*$.
\end{proposition}
\begin{proof}
By the minimax theorem, $v^*=\min_{y'}\max_a(My')_a\le\max_a(My)_a$. For the
second claim take matching pennies, $M=\bigl(\begin{smallmatrix}1&-1\\-1&1\end{smallmatrix}\bigr)$,
with $v^*=0$ and $y=(\tfrac12,\tfrac12)$. Every pure action is a best response,
and each loses 1 to some reply.
\end{proof}

\intuition{Against the \emph{modelled} opponent the pure choice is never worse
than the game's value. Against an opponent who anticipates it, a deterministic
best response can be exploited, whereas $\bar x$ could not be. The design bets
that ladder opponents neither compute this matrix nor exploit its argmax. The
level-1 component also assumes they reason one step about our options. These
are modelling choices, and the live evidence in Section~\ref{sec:evidence} is
too small to test them separately.}

\subsection{Determinization and strategy fusion}
\label{sec:fusion}
Sampling worlds and solving each with full knowledge is \emph{perfect-information
Monte Carlo} (PIMC). It is known to suffer from \emph{strategy fusion}, which
means acting as if one could choose differently in worlds one cannot tell
apart, and from \emph{non-locality} \citep{frankbasin1998}. Yet PIMC often
plays well in practice \citep{long2010}. Information-set search avoids fusion
at a higher cost \citep{cowling2012}. Here the turn decision avoids fusion at
the root: Equation~\eqref{eq:decision} commits to one $a$ for all worlds.
Fusion does enter forced-replacement scoring, which takes
$\frac1K\sum_w\max_aU_w(a)$.

\begin{proposition}[Forced-replacement scores are optimistic]
\label{prop:fusion}
For any values $U_w(a)$,
\[
  \frac1K\sum_w\max_aU_w(a)\ \ge\ \max_a\frac1K\sum_wU_w(a),
\]
with equality if and only if some single action maximizes $U_w$ in every world.
\end{proposition}
\begin{proof}
For each $a'$ and $w$, $\max_aU_w(a)\ge U_w(a')$. Averaging over $w$ and then
maximizing over $a'$ gives the inequality. If $a'$ maximizes in every world,
the two sides coincide. Conversely, if no action does, then for every $a'$
some world has a strict gap, so the average on the left exceeds every
right-hand average.
\end{proof}

\intuition{A replacement whose best follow-up depends on the hidden world is
scored as if we will learn the world before moving. That inflates its value,
so the rule favours flexible-looking replacements whose flexibility is not
real. Within each world both simulated sides also know each other's sets,
even when sheets were not shared, so the opponent model is omniscient about
us.}

\subsection{Monte Carlo error}
Each cell uses one stochastic sample of accuracy, damage rolls and critical
hits. Each cell has its own seed, so two of our actions are not compared under
common random numbers. With $K=6$ worlds, the averaged difference between two
actions has world-sampling noise plus within-cell chance noise.

\begin{proposition}[Ranking two actions]
\label{prop:mc}
Suppose the per-world differences $D_w=U_w(a)-U_w(a')$ are independent, lie in
an interval of width $2\Delta$, and have mean $\gamma>0$. Then the average
over $K$ worlds ranks $a'$ above $a$ with probability at most
$\exp(-K\gamma^2/(2\Delta^2))$.
\end{proposition}
\begin{proof}
This is Hoeffding's inequality applied to the event $\overline D\le0$ with
deviation $\gamma$ \citep{hoeffding}.
\end{proof}

With $K=6$, the bound is useful only when the gap $\gamma$ is a large fraction
of $\Delta$. Small value differences are effectively ranked at random. Common
random numbers across rows, more seeds per cell, or adaptive allocation of
worlds to close decisions would all reduce this noise.

\subsection{Preview, Mega and the learned value}
Team preview is chosen by the actor, greedily. A value-based preview mode
(score 90 plans by the position after preview, averaged over 6 worlds and 3
sampled opposing bring-fours) exists but is untested. Mega Evolution is always
taken when available, for us and for the opponent model, so the agent never
considers saving it. A small value network can replace $V$: 14 features per
Pok\'{e}mon plus field and matchup features, a $128\to128\to1$ ReLU MLP,
trained by binary cross-entropy on omniscient self-play. It blends in as
$(1-\beta)V+\beta\cdot8\tanh(z/2)$, but the deployment uses $\beta=0$, and the
network is neither trained to completion nor gated.

\source{\code{search/agent.py} (worlds, pruning, move-id rewriting, decision
rule, forced replacements), \code{search/engine.cjs} (rebuild, enumeration,
screening, payoffs, regret matching, evaluation), \code{search/priors.py},
\code{search/tracker.py}, \code{search/valuefeat.cjs}, \code{search/train_value.py};
\code{docs/search-agent.md}.}
